% GENERATED FILE. Edit canonical lessons, then run npm run generate:books.
% canonical-source-hash: fnv1a64:069e6063a2b107b4
% canonical-lessons: RU-C433-shum, RU-C433-meshat, RU-C433-tarakan, RU-C433-seyf, RU-C433-nepoladka

\chapter{Noisy neighbours}
\label{ch:ru-noisy-neighbours}
\hlchaptermodality{433}

\begin{chapteropening}
\textbf{By the end of this chapter:} \emph{I can complain about a room: noise, being disturbed, cockroaches, the safe and a fault.}

Complete the last lesson of chapter 433: I can complain about a room: noise, being disturbed, cockroaches, the safe and a fault.
\end{chapteropening}

\section[\texorpdfstring{\ru{шум} (shum)}{шум (shum)}]{\ru{шум} (shum) — noise}

\label{lesson:RU-C433-shum}



\begin{quote}
Before the new one: say the Russian for to flow, then the Russian for a flood.
\end{quote}

\subsection*{You'll want to know: \ru{шум}}
\textbf{\ru{шум}} (\emph{shum}) — \textquotedblleft{}noise\textquotedblright{}.

\begin{quote}
\ru{Из}-\ru{за} \ru{шума} \ru{на} \ru{улице} \ru{я} \ru{плохо} \ru{спал}.
\end{quote}

\begin{quote}
\emph{Iz-za shúma na úlitse ya plókho spal.}
\end{quote}

\begin{quote}
Because of the noise in the street I slept badly.
\end{quote}

Masculine. \emph{\ru{Шуметь}} — to make noise.

\begin{cousinweb}
\emph{\ru{шум}} is inherited from Proto-Slavic \emph{šumŭ} (EtymDB), the root of \emph{\ru{шумный}} 'noisy'.
\end{cousinweb}

\begin{grammarlens}[title={What you've built}]
One more word for telling, explaining and describing at B1.
\end{grammarlens}

\subsection*{Your turn}
\begin{itemize}
\raggedright
  \item \emph{Say it:} \emph{shum}
  \item \emph{Say it:} \emph{shum}, once more
  \item \emph{Say it:} a sentence of your own with \emph{shum}
  \item \emph{Recall:} say the Russian for to flow, then the Russian for a flood, then say \emph{shum} again
\end{itemize}

\begin{tcolorbox}[breakable,colback=teal!4,colframe=teal!35!black,title={Before you move on}]
What does \ru{шум} mean? (\textquotedblleft{}Noise\textquotedblright{}.) Say it once more.
\end{tcolorbox}
\section[\texorpdfstring{\ru{мешать} (meshát)}{мешать (meshát)}]{\ru{мешать} (meshát) — to disturb, to bother; to stir}

\label{lesson:RU-C433-meshat}



\begin{quote}
Before the new one: say the Russian for a flood, then the Russian for noise.
\end{quote}

\subsection*{You'll want to know: \ru{мешать}}
\textbf{\ru{мешать}} (\emph{meshát}) — \textquotedblleft{}to disturb, to bother; to stir\textquotedblright{}.

\begin{quote}
\ru{Музыка} \ru{соседей} \ru{мешает} \ru{нам} \ru{спать}.
\end{quote}

\begin{quote}
\emph{Múzyka sosédey mesháyet nam spat.}
\end{quote}

\begin{quote}
The neighbours' music keeps us from sleeping.
\end{quote}

Imperfective; perfective \emph{\ru{помешать}}. Takes the dative: \emph{\ru{мешать} \ru{соседям}}. Also to stir: \emph{\ru{мешать} \ru{суп}}.

\begin{cousinweb}
\emph{\ru{мешать}} is usually traced to an old Slavic verb meaning 'to mix', which fits \emph{\ru{мешать} \ru{суп}} 'to stir soup'; to disturb is to mix things up for someone. EtymDB lists the word without a source.
\end{cousinweb}

\begin{grammarlens}[title={What you've built}]
One more word for telling, explaining and describing at B1.
\end{grammarlens}

\subsection*{Your turn}
\begin{itemize}
\raggedright
  \item \emph{Say it:} \emph{meshát}
  \item \emph{Say it:} \emph{meshát}, once more
  \item \emph{Say it:} a sentence of your own with \emph{meshát}
  \item \emph{Recall:} say the Russian for a flood, then the Russian for noise, then say \emph{meshát} again
\end{itemize}

\begin{tcolorbox}[breakable,colback=teal!4,colframe=teal!35!black,title={Before you move on}]
What does \ru{мешать} mean? (\textquotedblleft{}To disturb, to bother; to stir\textquotedblright{}.) Say it once more.
\end{tcolorbox}
\section[\texorpdfstring{\ru{таракан} (tarakán)}{таракан (tarakán)}]{\ru{таракан} (tarakán) — a cockroach}

\label{lesson:RU-C433-tarakan}



\begin{quote}
Before the new one: say the Russian for noise, then the Russian for to disturb.
\end{quote}

\subsection*{You'll want to know: \ru{таракан}}
\textbf{\ru{таракан}} (\emph{tarakán}) — \textquotedblleft{}a cockroach\textquotedblright{}.

\begin{quote}
\ru{На} \ru{кухне} \ru{мы} \ru{увидели} \ru{таракана}.
\end{quote}

\begin{quote}
\emph{Na kúkhne my uvídeli tarakána.}
\end{quote}

\begin{quote}
We saw a cockroach in the kitchen.
\end{quote}

Masculine. Colloquial: \emph{\ru{тараканы} \ru{в} \ru{голове}} — odd ideas.

\begin{cousinweb}
The origin of \emph{\ru{таракан}} is not settled; it is often traced to a Turkic word. EtymDB lists the Russian word without a source.
\end{cousinweb}

\begin{grammarlens}[title={What you've built}]
One more word for telling, explaining and describing at B1.
\end{grammarlens}

\subsection*{Your turn}
\begin{itemize}
\raggedright
  \item \emph{Say it:} \emph{tarakán}
  \item \emph{Say it:} \emph{tarakán}, once more
  \item \emph{Say it:} a sentence of your own with \emph{tarakán}
  \item \emph{Recall:} say the Russian for noise, then the Russian for to disturb, then say \emph{tarakán} again
\end{itemize}

\begin{tcolorbox}[breakable,colback=teal!4,colframe=teal!35!black,title={Before you move on}]
What does \ru{таракан} mean? (\textquotedblleft{}A cockroach\textquotedblright{}.) Say it once more.
\end{tcolorbox}
\section[\texorpdfstring{\ru{сейф} (seyf)}{сейф (seyf)}]{\ru{сейф} (seyf) — a safe}

\label{lesson:RU-C433-seyf}



\begin{quote}
Before the new one: say the Russian for to disturb, then the Russian for a cockroach.
\end{quote}

\subsection*{You'll want to know: \ru{сейф}}
\textbf{\ru{сейф}} (\emph{seyf}) — \textquotedblleft{}a safe\textquotedblright{}.

\begin{quote}
\ru{Паспорт} \ru{и} \ru{деньги} \ru{лежат} \ru{в} \ru{сейфе}.
\end{quote}

\begin{quote}
\emph{Pásport i déngi lezhát v séyfe.}
\end{quote}

\begin{quote}
The passport and the money are in the safe.
\end{quote}

Masculine.

\begin{cousinweb}
\emph{\ru{сейф}} was borrowed from English \emph{safe} (EtymDB), a strong box.
\end{cousinweb}

\begin{grammarlens}[title={What you've built}]
One more word for telling, explaining and describing at B1.
\end{grammarlens}

\subsection*{Your turn}
\begin{itemize}
\raggedright
  \item \emph{Say it:} \emph{seyf}
  \item \emph{Say it:} \emph{seyf}, once more
  \item \emph{Say it:} a sentence of your own with \emph{seyf}
  \item \emph{Recall:} say the Russian for to disturb, then the Russian for a cockroach, then say \emph{seyf} again
\end{itemize}

\begin{tcolorbox}[breakable,colback=teal!4,colframe=teal!35!black,title={Before you move on}]
What does \ru{сейф} mean? (\textquotedblleft{}A safe\textquotedblright{}.) Say it once more.
\end{tcolorbox}
\section[\texorpdfstring{\ru{неполадка} (nepoládka)}{неполадка (nepoládka)}]{\ru{неполадка} (nepoládka) — a fault, a problem (with equipment)}

\label{lesson:RU-C433-nepoladka}



\begin{quote}
Before the new one: say the Russian for a cockroach, then the Russian for a safe.
\end{quote}

\subsection*{You'll want to know: \ru{неполадка}}
\textbf{\ru{неполадка}} (\emph{nepoládka}) — \textquotedblleft{}a fault, a problem (with equipment)\textquotedblright{}.

\begin{quote}
\ru{В} \ru{номере} \ru{была} \ru{неполадка} \ru{с} \ru{душем}.
\end{quote}

\begin{quote}
\emph{V nómere bylá nepoládka s dúshem.}
\end{quote}

\begin{quote}
There was a problem with the shower in the room.
\end{quote}

Feminine; often plural: \emph{\ru{неполадки}}.

\begin{cousinweb}
\emph{\ru{неполадка}} is built from \emph{\ru{не}-} 'not' and \emph{\ru{поладить}} 'to get on', on \emph{\ru{лад}} 'harmony, good order', which is inherited from Proto-Slavic \emph{ladŭ} (EtymDB).
\end{cousinweb}

\begin{grammarlens}[title={What you've built}]
One more word for telling, explaining and describing at B1.
\end{grammarlens}

\subsection*{Your turn}
\begin{itemize}
\raggedright
  \item \emph{Say it:} \emph{nepoládka}
  \item \emph{Say it:} \emph{nepoládka}, once more
  \item \emph{Say it:} a sentence of your own with \emph{nepoládka}
  \item \emph{Recall:} say the Russian for a cockroach, then the Russian for a safe, then say \emph{nepoládka} again
\end{itemize}

\begin{tcolorbox}[breakable,colback=teal!4,colframe=teal!35!black,title={Before you move on}]
What does \ru{неполадка} mean? (\textquotedblleft{}A fault, a problem (with equipment)\textquotedblright{}.) Say it once more.
\end{tcolorbox}
